% GENERATED FILE. Edit canonical lessons, then run npm run generate:books.
% canonical-source-hash: fnv1a64:1529fb391c677665
% canonical-lessons: ZH-C199-chima, ZH-C199-bao, ZH-C199-fei, ZH-C199-jin, ZH-C199-tiao-pick, ZH-R199-review-194-199

\chapter{Sizes}
\label{ch:zh-sizes}
\hlchaptermodality{199}

\begin{chapteropening}
\textbf{By the end of this chapter:} \emph{I can ask what size someone wears, say clothes are thin, baggy or tight, and pick one out.}

Complete the last lesson of chapter 199: I can ask what size someone wears, say clothes are thin, baggy or tight, and pick one out.
\end{chapteropening}

\section[chǐmǎ]{chǐmǎ — clothing size; shoe size}

\label{lesson:ZH-C199-chima}



\begin{quote}
Before the new one: say the Mandarin for engaged (of a phone line), then the Mandarin for phone signal.
\end{quote}

\subsection*{You'll want to know: chǐmǎ}
\emph{chǐmǎ} — \textquotedblleft{}clothing size; shoe size\textquotedblright{}.

\begin{quote}
\emph{nǐ chuān duōdà de chǐmǎ} — What size do you wear?
\end{quote}

Its characters come later; for now it is a word for the ear.

Two third tones meet in this word: chǐ is said with a rising tone before mǎ, and the spelling keeps both marks.

\begin{cousinweb}
\emph{chǐmǎ} is built from \emph{chǐ}, a measure of length, and \emph{mǎ} 'a number sign', both inherited from Middle Chinese, \emph{tsyhek} and \emph{maeX} (Guangyun, in Baxter's transcription). The dictionary explains the first as a measure of length and glosses the second as a stone like jade; how the pair came to mean the size of clothes is not settled.
\end{cousinweb}

\begin{grammarlens}[title={What you've built}]
One more word for asking for what I want.
\end{grammarlens}

\subsection*{Your turn}
\begin{itemize}
\raggedright
  \item \emph{Say it:} \emph{chǐmǎ}
  \item \emph{Say it:} \emph{chǐmǎ}, once more
  \item \emph{Say it:} \emph{nǐ chuān duōdà de chǐmǎ} — What size do you wear?
  \item \emph{Recall:} say the Mandarin for engaged (of a phone line), then the Mandarin for phone signal, then say \emph{chǐmǎ} again
\end{itemize}

\begin{tcolorbox}[breakable,colback=teal!4,colframe=teal!35!black,title={Before you move on}]
What does chǐmǎ mean? (\textquotedblleft{}Clothing size; shoe size\textquotedblright{}.) Say it once more.
\end{tcolorbox}
\section[báo]{báo — thin (of cloth or paper)}

\label{lesson:ZH-C199-bao}



\begin{quote}
Before the new one: say the Mandarin for phone signal, then the Mandarin for clothing size.
\end{quote}

\subsection*{You'll want to know: báo}
\emph{báo} — \textquotedblleft{}thin (of cloth or paper)\textquotedblright{}.

\begin{quote}
\emph{nǐ de yīfu tài báo le} — Your clothes are too thin.
\end{quote}

Its characters come later; for now it is a word for the ear.

\begin{cousinweb}
\emph{báo} is inherited from Middle Chinese \emph{bak} (Guangyun, in Baxter's transcription), glossed by the pair 'thick and thin'. Mandarin has two outcomes of that form: \emph{báo} in everyday speech, and \emph{bó} in bookish words.
\end{cousinweb}

\begin{grammarlens}[title={What you've built}]
One more word for asking for what I want.
\end{grammarlens}

\subsection*{Your turn}
\begin{itemize}
\raggedright
  \item \emph{Say it:} \emph{báo}
  \item \emph{Say it:} \emph{báo}, once more
  \item \emph{Say it:} \emph{nǐ de yīfu tài báo le} — Your clothes are too thin.
  \item \emph{Recall:} say the Mandarin for phone signal, then the Mandarin for clothing size, then say \emph{báo} again
\end{itemize}

\begin{tcolorbox}[breakable,colback=teal!4,colframe=teal!35!black,title={Before you move on}]
What does báo mean? (\textquotedblleft{}Thin (of cloth or paper)\textquotedblright{}.) Say it once more.
\end{tcolorbox}
\section[féi]{féi — loose; baggy (of clothes)}

\label{lesson:ZH-C199-fei}



\begin{quote}
Before the new one: say the Mandarin for clothing size, then the Mandarin for thin (of cloth or paper).
\end{quote}

\subsection*{You'll want to know: féi}
\emph{féi} — \textquotedblleft{}loose; baggy (of clothes)\textquotedblright{}.

\begin{quote}
\emph{zhè tiáo kùzi yǒudiǎnr féi} — These trousers are a bit baggy.
\end{quote}

Its characters come later; for now it is a word for the ear.

\begin{cousinweb}
\emph{féi} is inherited from Middle Chinese \emph{bj+j} (Guangyun, in Baxter's transcription), glossed 'fat' and 'fleshy'. Its use for loose clothes is not given there, so how it came is not settled.
\end{cousinweb}

\begin{grammarlens}[title={What you've built}]
One more word for asking for what I want.
\end{grammarlens}

\subsection*{Your turn}
\begin{itemize}
\raggedright
  \item \emph{Say it:} \emph{féi}
  \item \emph{Say it:} \emph{féi}, once more
  \item \emph{Say it:} \emph{zhè tiáo kùzi yǒudiǎnr féi} — These trousers are a bit baggy.
  \item \emph{Recall:} say the Mandarin for clothing size, then the Mandarin for thin (of cloth or paper), then say \emph{féi} again
\end{itemize}

\begin{tcolorbox}[breakable,colback=teal!4,colframe=teal!35!black,title={Before you move on}]
What does féi mean? (\textquotedblleft{}Loose; baggy (of clothes)\textquotedblright{}.) Say it once more.
\end{tcolorbox}
\section[jǐn]{jǐn — tight}

\label{lesson:ZH-C199-jin}



\begin{quote}
Before the new one: say the Mandarin for thin (of cloth or paper), then the Mandarin for baggy.
\end{quote}

\subsection*{You'll want to know: jǐn}
\emph{jǐn} — \textquotedblleft{}tight\textquotedblright{}.

\begin{quote}
\emph{zhè shuāng xié tài jǐn le} — These shoes are too tight.
\end{quote}

Its characters come later; for now it is a word for the ear.

\begin{cousinweb}
\emph{jǐn} is inherited from Middle Chinese \emph{kjinX} (Guangyun, in Baxter's transcription), glossed 'tightly bound'.
\end{cousinweb}

\begin{grammarlens}[title={What you've built}]
One more word for asking for what I want.
\end{grammarlens}

\subsection*{Your turn}
\begin{itemize}
\raggedright
  \item \emph{Say it:} \emph{jǐn}
  \item \emph{Say it:} \emph{jǐn}, once more
  \item \emph{Say it:} \emph{zhè shuāng xié tài jǐn le} — These shoes are too tight.
  \item \emph{Recall:} say the Mandarin for thin (of cloth or paper), then the Mandarin for baggy, then say \emph{jǐn} again
\end{itemize}

\begin{tcolorbox}[breakable,colback=teal!4,colframe=teal!35!black,title={Before you move on}]
What does jǐn mean? (\textquotedblleft{}Tight\textquotedblright{}.) Say it once more.
\end{tcolorbox}
\section[tiāo]{tiāo — to pick out; to choose}

\label{lesson:ZH-C199-tiao-pick}



\begin{quote}
Before the new one: say the Mandarin for baggy, then the Mandarin for tight.
\end{quote}

\subsection*{You'll want to know: tiāo}
\emph{tiāo} — \textquotedblleft{}to pick out; to choose\textquotedblright{}.

\begin{quote}
\emph{qǐng bāng wǒ tiāo yí ge chǐmǎ} — Could you help me pick out a size, please?
\end{quote}

Its characters come later; for now it is a word for the ear.

\begin{cousinweb}
\emph{tiāo} is inherited from Middle Chinese \emph{thew} (Guangyun, in Baxter's transcription), glossed 'to stir up'. The sense 'to pick out' is not given there, so how it came is not settled.
\end{cousinweb}

\begin{grammarlens}[title={What you've built}]
One more word for asking for what I want.
\end{grammarlens}

\subsection*{Your turn}
\begin{itemize}
\raggedright
  \item \emph{Say it:} \emph{tiāo}
  \item \emph{Say it:} \emph{tiāo}, once more
  \item \emph{Say it:} \emph{qǐng bāng wǒ tiāo yí ge chǐmǎ} — Could you help me pick out a size, please?
  \item \emph{Recall:} say the Mandarin for baggy, then the Mandarin for tight, then say \emph{tiāo} again
\end{itemize}

\begin{tcolorbox}[breakable,colback=teal!4,colframe=teal!35!black,title={Before you move on}]
What does tiāo mean? (\textquotedblleft{}To pick out; to choose\textquotedblright{}.) Say it once more.
\end{tcolorbox}
\section[(dialogue)]{Looking back: from Roads and Crossings to Sizes}

\label{lesson:ZH-R199-review-194-199}



\begin{quote}
Say the Mandarin for a street for traffic and to go on foot.
\end{quote}

\subsection*{Your turn}
\begin{itemize}
\raggedright
  \item \emph{Say it:} a street for traffic, to go on foot, zebra crossing, public square, traffic lights
  \item \emph{Say it:} to put money by, to take out money, credit card, password, bank account
  \item \emph{Say it:} train or bus ticket, plane ticket, scheduled flight, to take off, to reach a destination
  \item \emph{Say it:} to ask the way, to find, to lose your way, straight on, route
  \item \emph{Say it:} to switch off, to turn a phone on, to leave a message, engaged (of a phone line), phone signal
  \item \emph{Say it:} clothing size, thin (of cloth or paper), baggy, tight, to pick out
  \item \emph{Read it:} \zh{路} \zh{票} \zh{找} \zh{关} — each one aloud, then a word that holds it
\end{itemize}

\begin{tcolorbox}[breakable,colback=teal!4,colframe=teal!35!black,title={Before you move on}]
A street for traffic? (\textbf{\zh{马路}}, \emph{mǎlù}.) To go on foot? (\emph{zǒulù}.) Zebra crossing? (\emph{bānmǎxiàn}.) Public square? (\emph{guǎngchǎng}.) Traffic lights? (\emph{hónglǜdēng}.) To put money by? (\emph{cúnqián}.) To take out money? (\emph{qǔqián}.) Credit card? (\emph{xìnyòngkǎ}.) Password? (\emph{mìmǎ}.) Bank account? (\emph{zhànghù}.) Train or bus ticket? (\textbf{\zh{车票}}, \emph{chēpiào}.) Plane ticket? (\textbf{\zh{机票}}, \emph{jīpiào}.) Scheduled flight? (\emph{hángbān}.) To take off? (\emph{qǐfēi}.) To reach a destination? (\emph{dàodá}.) To ask the way? (\textbf{\zh{问路}}, \emph{wènlù}.) To find? (\textbf{\zh{找到}}, \emph{zhǎodào}.) To lose your way? (\emph{mílù}.) Straight on? (\emph{yìzhí}.) Route? (\emph{lùxiàn}.) To switch off? (\textbf{\zh{关机}}, \emph{guānjī}.) To turn a phone on? (\textbf{\zh{开机}}, \emph{kāijī}.) To leave a message? (\emph{liúyán}.) Engaged (of a phone line)? (\emph{zhànxiàn}.) Phone signal? (\emph{xìnhào}.) Clothing size? (\emph{chǐmǎ}.) Thin (of cloth or paper)? (\emph{báo}.) Baggy? (\emph{féi}.) Tight? (\emph{jǐn}.) To pick out? (\emph{tiāo}.)
\end{tcolorbox}
